\documentclass[10pt,a4paper]{amsart}
\usepackage[margin=19mm]{geometry}
\usepackage{amsmath,amssymb,mathtools}
\usepackage[scr=rsfs,cal=euler]{mathalfa}
\usepackage{xfrac}
\usepackage[unicode,hidelinks,bookmarks=false]{hyperref}
\newcommand{\ie}{i.e.\ }
\newcommand{\cf}{cf.\ }
\newcommand{\resp}{respectively\ }
\newcommand{\Subsection}{\subsection}
\providecommand{\nobreakdash}{\nobreak\hbox{-}}
\setlength{\emergencystretch}{2em}
\multlinegap0pt
\usepackage{xcolor}
\usepackage{amssymb}
\usepackage{tikz}
\newtheorem{proposition}{Proposition}
\title{Simultaneous extensions of states on quantum logics}
\author{Dominika Bure\v{s}ov\'{a}, Pavel Pt\'{a}k, Jan \v{S}evic}
\date{Source: arXiv:2609.10770v1 (CC BY 4.0)}
\begin{document}
\maketitle

\noindent\emph{Source and licence.} Simultaneous extensions of states on quantum logics. Version arXiv:2609.10770v1 (CC BY 4.0), distributed by arXiv under the Creative Commons Attribution 4.0 International licence, http://creativecommons.org/licenses/by/4.0/, as recorded in the arXiv licence metadata for this exact version and shown on the abstract page https://arxiv.org/abs/2609.10770v1. Full licence notice: \texttt{licenses/arxiv-2609.10770-CC-BY-4.0.txt}.\par\smallskip

\noindent\textit{Editorial scope.} Counted unit: the article's proposition Prop2 (source0.tex lines 302--304). The article's complete original proof of it is delivered with it (source0.tex lines 306--318, one \texttt{proof} environment of 13 lines). No mathematics is written, repaired or re-derived: the statement and the proof are the article's own lines, in the article's own notation, and no retained line is substituted or re-flowed.  No further source-local context is needed: every cross-reference of every retained line resolves inside the two delivered spans.  Nothing was omitted from any retained span, so the package carries no omission record.  No retained line contains a citation, so the wrapper carries no bibliography and no bibliography entry.  No retained line contains a figure environment, an image-inclusion command or drawing code, so no image is delivered and the package declares no asset dependency.  Editorial formatting only, in addition: an AMS \texttt{amsart} preamble that restates the article's own declarations and macros used by the retained lines, namely the environments \texttt{proposition} on separate counters, together with the standard packages those lines need.  The retained environments print in this document as Proposition 1 in order of appearance, whereas the article prints the same items with the numbering of its own full document.  The complete native source of the article as distributed in the arXiv e-print for this exact version is bundled unchanged as \texttt{sources/209948/source0.tex}.
\begin{proposition}\label{Prop2}
Let $A$ be a Boolean subalgebra of a logic $L$. Let $L$ be $0$--$1$ distinguishing. Then any state $s\in \mathcal{S}(A)$ can be extended over $L$ as a~state.
\end{proposition}

\begin{proof}
Let $\mathcal F$ be the collection of all finite subalgebras of $A$. If
$F=\{0,1,a_1,\dots,a_n\}\in\mathcal F$,
then there is a state $v_F\in \mathcal{S}(L)$ such that $v_F$ restricted to $F$ equals to $s$.
Indeed, for each $a_i$ there is a state $t_i$ on $L$ with
$t_i(a_i)=1$ and we set $v_F=\sum_{i=1}^{n}s(a_i)t_i$.
Since $\sum_{i \leq n}s(a_i)=1$, we see that $v_F\in \mathcal{S}(L)$.
%Then $v_F\in \mathcal{S}(L)$.\mathcal{S}(L)

The collection of finite subalgebras of $A$ is a directed set under inclusion, so the compactness of $\mathcal{S}(L)$ gives us a state
$t$ that we obtain as a~directed limit $t=\lim\limits_{F\in\mathcal F}v_F$.
Obviously, $t$ extends $s$.
\end{proof}

\end{document}
